\chapter{Fundamentação Teórica}
\label{cap:rev.teorica}
%\thispagestyle{plain}
\graphicspath{{./Cap2_Revisao_Teorica/Figures/}}

A evolução dos sistemas de controle automático reflete o desenvolvimento tecnológico da humanidade, desde os primeiros mecanismos de realimentação até os modernos sistemas baseados em inteligência artificial. Os primórdios do controle automático remontam aos reguladores mecânicos do século XVIII, como o regulador centrífugo de James Watt para máquinas a vapor, estabelecendo os fundamentos do controle por realimentação. A formalização matemática destes conceitos ocorreu no século XX, com o desenvolvimento da teoria de controle clássico, caracterizada pela utilização de técnicas baseadas em função de transferência e análise no domínio da frequência \cite{ogata2009}.

O controlador PID (Proporcional-Integral-Derivativo) emergiu como a solução mais amplamente adotada no controle clássico, oferecendo simplicidade de implementação e eficácia em uma ampla gama de aplicações industriais. Sua popularidade deve-se à capacidade de fornecer controle estável e responsivo através do ajuste de três parâmetros fundamentais: ganho proporcional, tempo integral e tempo derivativo \cite{villaca2013}.

A década de 1960 marcou o início da era do controle moderno, caracterizada pela representação de sistemas em espaço de estados e pela utilização de métodos de otimização. Esta abordagem permitiu o tratamento de sistemas multivariáveis e não lineares de forma mais sistemática, superando limitações do controle clássico. O controle moderno introduziu conceitos como controlabilidade, observabilidade e estabilidade no sentido de Lyapunov, com técnicas como o regulador linear quadrático (LQR) e o filtro de Kalman proporcionando ferramentas poderosas para o projeto de controladores ótimos \cite{franklin2015}.

O reconhecimento de que os modelos matemáticos são sempre aproximações da realidade levou ao desenvolvimento do controle robusto nas décadas de 1980 e 1990. Esta abordagem considera explicitamente as incertezas do modelo, garantindo estabilidade e desempenho satisfatório mesmo na presença de variações paramétricas e distúrbios não modelados.

A lógica Fuzzy, introduzida por Lotfi Zadeh em 1965, representa uma abordagem fundamentalmente diferente para o controle de sistemas, baseada na capacidade humana de raciocinar com informações imprecisas e tomar decisões em ambientes de incerteza. Esta técnica emerge como uma ponte entre o controle convencional e a inteligência artificial aplicada, posicionando-se como uma alternativa às limitações dos métodos clássicos e modernos \cite{zadeh1965}.

O desenvolvimento do controle Fuzzy surge em resposta às dificuldades encontradas pelos métodos convencionais em lidar com sistemas complexos, não lineares e mal definidos matematicamente. Enquanto o controle clássico e moderno dependem de modelos matemáticos precisos, o controle Fuzzy permite a incorporação de conhecimento heurístico e experiência operacional diretamente no algoritmo de controle, dispensando a necessidade de modelagem matemática rigorosa \cite{mamdani1975}.

O controle Fuzzy posiciona-se na evolução dos sistemas de controle como uma abordagem complementar que preenche lacunas deixadas pelos métodos convencionais. Sua principal contribuição reside na capacidade de tratar incertezas e imprecisões de forma natural, oferecendo robustez e simplicidade de implementação em aplicações onde modelos matemáticos precisos são difíceis de obter ou onde existe conhecimento heurístico valioso disponível \cite{ross2010}.

A integração eficiente de sistemas em ambientes industriais depende fundamentalmente da capacidade de comunicação entre dispositivos heterogêneos. A evolução dos protocolos de comunicação industrial reflete a crescente necessidade de interoperabilidade, segurança e eficiência na troca de informações. A comunicação industrial evoluiu de sistemas proprietários isolados para padrões abertos e interoperáveis, onde inicialmente cada fabricante desenvolvia seus próprios protocolos, resultando em "ilhas de automação" com limitada capacidade de integração \cite{stenerson2015}.

A necessidade de interoperabilidade levou ao desenvolvimento de padrões como Modbus, Profibus e DeviceNet. O OPC (Open Platform Communications) foi desenvolvido para resolver problemas de interoperabilidade em automação industrial, estabelecendo um padrão para comunicação entre aplicações de diferentes fornecedores. Baseado inicialmente na tecnologia DCOM da Microsoft, o OPC Classic proporcionou avanços significativos, mas suas limitações, incluindo dependência de plataforma e problemas de segurança, levaram ao desenvolvimento do OPC UA (Unified Architecture) \cite{mahnke2009}.

\cite{carvalho2023} investigaram a digitalização de uma planta industrial utilizando o protocolo de comunicação OPC-UA, demonstrando como a Indústria 4.0 está impulsionando transformações profundas na indústria através de soluções que aprimoram a produção. A pesquisa destaca que a utilização do protocolo OPC UA para conectar módulos ciberfísicos e dispositivos à nuvem possibilita análise de dados em tempo real para identificar status atual, problemas potenciais e tomar decisões mais eficazes.

\cite{souza2024} conduziu um estudo de caso sobre controle e supervisão de processos utilizando o padrão OPC, focando na superação dos desafios de integração entre equipamentos e softwares de diferentes fabricantes devido a protocolos proprietários. O trabalho utilizou o CLP Codesys e o ambiente simulado do Factory IO, demonstrando como o padrão OPC UA promove a troca de dados e comandos entre sistemas distintos.

O OPC UA representa uma arquitetura completamente nova, projetada para ser independente de plataforma, segura, escalável e orientada a serviços. Utiliza uma arquitetura em camadas que separa a lógica de aplicação dos detalhes de comunicação, com modelo de informação baseado em uma estrutura hierárquica de nós e referências que permite representar informações complexas de forma padronizada e extensível \cite{leitner2006}.

A utilização de protocolos de comunicação padronizados em plantas industriais transcende aspectos puramente técnicos, representando um facilitador fundamental para a implementação da Indústria 4.0. O OPC UA alinha-se perfeitamente com os princípios da Indústria 4.0, facilitando a integração de sistemas ciberfísicos, permitindo comunicação entre sistemas heterogêneos, suportando diversos paradigmas de comunicação e adaptando-se desde dispositivos de campo até sistemas na nuvem \cite{jasperneite2012}.

A combinação de protocolos padronizados como OPC UA com técnicas de controle inteligente, incluindo lógica Fuzzy, oferece oportunidades únicas para a criação de sistemas de automação mais eficientes e adaptáveis. Esta integração permite distribuição de inteligência em diferentes níveis da hierarquia de controle, flexibilidade arquitetural onde mudanças nos algoritmos de controle não afetam a infraestrutura de comunicação, reutilização de componentes e manutenibilidade através da separação clara entre lógica de controle e comunicação.

\cite{oliveira2023} desenvolveu um estudo sobre gerenciamento de nível em reservatório de líquidos utilizando lógica Fuzzy e controle PID, implementando uma arquitetura integrada baseada em OPC para comunicação entre sistemas. O trabalho demonstrou controle satisfatório do sistema simulado de tanque, ilustrando como a combinação de técnicas de controle inteligente com protocolos de comunicação padronizados pode resultar em sistemas de controle mais eficazes e facilmente integráveis.

A convergência entre controle inteligente e comunicação padronizada representa um paradigma fundamental para a automação industrial moderna, proporcionando a base tecnológica necessária para implementar os conceitos da Indústria 4.0 de forma eficaz e sustentável. Esta integração facilita a implementação de estratégias de manutenção preditiva e proativa, contribui para a democratização do acesso a tecnologias avançadas e permite que empresas de diferentes portes implementem soluções sofisticadas de automação sem dependência de fornecedores específicos.

